% 2023 新高考Ⅱ卷 19：患病者、未患病者两组指标的频率分布直方图（左右并排）
% 绘制依据：源 PDF 原图。组距均为 5，纵轴「频率/组距」，横轴末端标「指标」。
% 左（患病者）刻度 0,95,100,…,130，七组高度 0.002/0.012/0.034/0.036/0.040/0.040/0.036；
% 右（未患病者）刻度 0,70,75,…,105，七组高度 0.038/0.040/0.040/0.036/0.034/0.010/0.002。
% （两组 Σ高度×5 均为 1。）横轴 0 与首刻度之间、纵轴 0 与最低标注之间均画断轴锯齿；
% 每个纵轴标注值处引水平虚线至对应柱顶。
% 版式：纵轴按 h(v)=0.45+(v−0.002)×125 (cm) 线性映射（故 0.002 的间隔恒为 0.25 cm），
% 柱宽 0.52 cm，右图 xshift=5.9 cm，整图宽约 11.9 cm；密集刻度用 \tiny 以免叠字。
\begin{tikzpicture}[font=\small]
  % ================= 左：患病者 =================
  \begin{scope}
    % 断轴的横轴与纵轴
    \draw (0,0) -- (0.10,0);
    \draw (0.10,0) -- (0.16,0.10) -- (0.24,-0.10) -- (0.30,0);
    \draw[->] (0.30,0) -- (4.55,0);
    \draw (0,0) -- (0,0.10);
    \draw (0,0.10) -- (0.09,0.16) -- (-0.09,0.29) -- (0,0.35);
    \draw[->] (0,0.35) -- (0,5.75);
    \node[anchor=north east, inner sep=1.5pt, font=\tiny] at (0,0) {0};
    \node[anchor=north, inner sep=2pt] at (4.85,0) {指标};
    \node[anchor=south, inner sep=3pt]
      at (0,5.75) {$\displaystyle\frac{\mbox{频率}}{\mbox{组距}}$};
    % 纵轴标注 + 水平虚线引至对应柱顶
    \foreach \h/\lab/\xe in {0.45/0.002/0.45, 1.70/0.012/0.97, 4.45/0.034/1.49,
                             4.70/0.036/2.01, 5.20/0.040/2.53} {
      \draw[dashed] (0,\h) -- (\xe,\h);
      \node[anchor=east, inner sep=2pt, font=\tiny] at (0,\h) {\lab};
    }
    % 七个矩形
    \foreach \i/\h in {0/0.45, 1/1.70, 2/4.45, 3/4.70, 4/5.20, 5/5.20, 6/4.70}
      \draw ({0.45+\i*0.52},0) rectangle ({0.45+(\i+1)*0.52},\h);
    % 横轴刻度
    \foreach \i/\lab in {0/95, 1/100, 2/105, 3/110, 4/115, 5/120, 6/125, 7/130} {
      \draw ({0.45+\i*0.52},0) -- ({0.45+\i*0.52},-0.07);
      \node[anchor=north, inner sep=1.5pt, font=\tiny] at ({0.45+\i*0.52},-0.07) {\lab};
    }
    \node at (2.27,-0.95) {患病者};
  \end{scope}
  % ================= 右：未患病者 =================
  \begin{scope}[xshift=5.9cm]
    \draw (0,0) -- (0.10,0);
    \draw (0.10,0) -- (0.16,0.10) -- (0.24,-0.10) -- (0.30,0);
    \draw[->] (0.30,0) -- (4.55,0);
    \draw (0,0) -- (0,0.10);
    \draw (0,0.10) -- (0.09,0.16) -- (-0.09,0.29) -- (0,0.35);
    \draw[->] (0,0.35) -- (0,5.75);
    \node[anchor=north east, inner sep=1.5pt, font=\tiny] at (0,0) {0};
    \node[anchor=north, inner sep=2pt] at (4.85,0) {指标};
    \node[anchor=south, inner sep=3pt]
      at (0,5.75) {$\displaystyle\frac{\mbox{频率}}{\mbox{组距}}$};
    \foreach \h/\lab/\xe in {0.45/0.002/3.57, 1.45/0.010/3.05, 4.45/0.034/2.53,
                             4.70/0.036/2.01, 4.95/0.038/0.45, 5.20/0.040/0.97} {
      \draw[dashed] (0,\h) -- (\xe,\h);
      \node[anchor=east, inner sep=2pt, font=\tiny] at (0,\h) {\lab};
    }
    \foreach \i/\h in {0/4.95, 1/5.20, 2/5.20, 3/4.70, 4/4.45, 5/1.45, 6/0.45}
      \draw ({0.45+\i*0.52},0) rectangle ({0.45+(\i+1)*0.52},\h);
    \foreach \i/\lab in {0/70, 1/75, 2/80, 3/85, 4/90, 5/95, 6/100, 7/105} {
      \draw ({0.45+\i*0.52},0) -- ({0.45+\i*0.52},-0.07);
      \node[anchor=north, inner sep=1.5pt, font=\tiny] at ({0.45+\i*0.52},-0.07) {\lab};
    }
    \node at (2.27,-0.95) {未患病者};
  \end{scope}
\end{tikzpicture}
